\chapter{Streptococcal infection}
\index{Streptococcal infection}

\medskip
\noindent\textit{Educational reference; always seek professional advice for individual care.}

\section*{Definition and Overview}
Streptococcal infections are illnesses caused by bacteria of the genus \textit{Streptococcus}. These are among the commonest bacteria to affect humans and produce an enormous range of disease, from trivial skin complaints to life-threatening sepsis. The two groups of greatest importance are group A streptococcus (\textit{Streptococcus pyogenes}), which causes sore throats, impetigo, cellulitis, and scarlet fever and is responsible for serious immune complications such as rheumatic fever, and group B streptococcus (\textit{Streptococcus agalactiae}), a cause of serious infection in newborns and of illness in older or unwell adults. Other streptococci live harmlessly in the mouth, bowel, and skin and only cause disease when they reach the wrong place, such as the blood or heart valves. Because the word "strep" is used loosely, this chapter gives an overview of the family of streptococcal infections: the common sites they affect, how they are transmitted, how they are diagnosed and treated, and when they become dangerous. Understanding the range is useful because the same bacterium can be harmless in one person and severe in another, and because prompt treatment prevents both spread and complications.

\section*{Epidemiology}
Streptococcal infections are among the most frequent bacterial illnesses worldwide. Group A streptococcus causes an estimated several hundred million throat and skin infections each year across the globe, and most children will experience at least one. Sore throats and school sores (impetigo) peak in school-aged children, in late winter and spring in temperate regions such as southern some countries, and in the wet season in the tropics. Group A streptococcus is also the leading infectious cause of preventable heart disease worldwide through rheumatic fever, and while this is now rare in high-income countries, the affected region still records significant numbers of cases among Aboriginal and Torres Strait Islander people, with some of the highest rates in the world in remote communities. Group B streptococcus colonises the vagina and bowel of roughly one in five healthy pregnant women, and while most babies born to carriers are healthy, it is the leading bacterial cause of serious illness in newborns, which is why pregnant women in many countries are offered routine screening. Invasive disease, where streptococci enter the bloodstream or deep tissues, is uncommon but serious, and its rates rise with age and with chronic illness.

\section*{Aetiology and Pathophysiology}
Streptococci are spherical bacteria that grow in chains and are classified by the antigens on their surface into groups A to G. They are spread by respiratory droplets, direct skin contact, or, for group B, from mother to baby during birth. The bacteria cause disease partly by direct tissue damage and partly by toxins and by immune responses that can attack the body's own tissues. The same strain can behave very differently in different people, and the site of infection determines the illness.

\begin{itemize}
    \item \textbf{Group A streptococcus:} causes throat infection (strep throat), impetigo or school sores, cellulitis, scarlet fever, and, rarely, invasive disease including toxic shock and flesh-eating infection
    \item \textbf{Group B streptococcus:} carried in the genital and bowel tract; the leading cause of early-onset sepsis and meningitis in newborns, and a cause of infection in older adults and people with diabetes
    \item \textbf{Pneumococcus (\textit{Streptococcus pneumoniae}):} a cause of pneumonia, sinusitis, middle-ear infection, and meningitis; prevented by the pneumococcal vaccine
    \item \textbf{Viridans group streptococci:} normal mouth residents that can cause endocarditis (infection of heart valves) if they reach the bloodstream, especially after dental work
    \item \textbf{Enterococcus:} a related bacterium that causes urinary and wound infections, particularly in hospitalised patients
    \item \textbf{Rheumatic fever:} an autoimmune reaction to group A streptococcus that can permanently damage heart valves
    \item \textbf{Post-streptococcal glomerulonephritis:} immune-mediated kidney inflammation following throat or skin infection
    \item \textbf{Antibiotic resistance:} increasing resistance complicates treatment, particularly among bacteria acquired in hospitals
\end{itemize}

\section*{Clinical Features}
The presentation depends on the site and type of infection.

\begin{itemize}
    \item \textbf{Sore throat and tonsillitis:} sudden severe throat pain, fever, pus on the tonsils, and swollen neck glands
    \item \textbf{Impetigo (school sores):} itchy red sores, usually around the nose and mouth, that weep and form golden crusts
    \item \textbf{Cellulitis:} a spreading, painful, hot red area of skin with fever, usually on a limb
    \item \textbf{Scarlet fever:} a fine, red, sandpaper-like rash, flushed cheeks, and a strawberry tongue, accompanying a strep throat
    \item \textbf{Newborn sepsis:} lethargy, poor feeding, fever or low temperature, and rapid breathing in the first week of life
    \item \textbf{Pneumonia:} cough, fever, breathlessness, and chest pain
    \item \textbf{Meningitis:} severe headache, stiff neck, fever, sensitivity to light, and altered consciousness
    \item \textbf{Invasive disease:} severe pain out of proportion to the skin findings, confusion, low blood pressure, and multi-organ failure
\end{itemize}

\section*{Differential Diagnosis}
Many other organisms cause similar pictures, and the diagnosis is often confirmed by testing.

\begin{itemize}
    \item \textbf{Viral sore throat and tonsillitis:} the commonest mimic; gradual onset with cough and runny nose, and no pus or response to antibiotics
    \item \textbf{Glandular fever:} severe throat pain and fatigue from Epstein-Barr virus, common in young adults; amoxicillin given at this time causes a typical rash
    \item \textbf{Staphylococcal skin infections:} staphylococci cause impetigo, boils, and cellulitis indistinguishable by eye, and are distinguished by culture
    \item \textbf{Viral exanthems:} measles, rubella, and other viruses produce rashes that can resemble scarlet fever
    \item \textbf{Other causes of neonatal sepsis:} \textit{E.~coli} and other organisms cause similar illness in newborns, which is why screening and prompt treatment matter
    \item \textbf{Other pneumonia and meningitis bacteria:} including \textit{Haemophilus}, meningococcus, and \textit{Listeria}, distinguished by cultures and antigen tests
\end{itemize}

\section*{When to Seek Medical Care}
Most streptococcal infections, such as a sore throat or minor impetigo, are assessed and treated by a general practitioner. A doctor should be consulted for a severe sore throat, a spreading red area of skin, a fever that does not settle, or a new rash, and for any illness in a newborn, a pregnant woman, or a person with a weakened immune system or chronic illness. People who carry group B streptococcus during pregnancy need planned intrapartum antibiotics, and any symptoms in the newborn should be assessed urgently. Because untreated group A infections can trigger rheumatic fever, throat and skin infections in high-risk communities should be treated promptly.

\textbf{Contact your local emergency services for:}
\begin{itemize}
    \item Difficulty breathing, or a blue or very pale complexion
    \item Confusion, severe drowsiness, or difficulty waking
    \item A rapidly spreading rash, bruising, or discoloured skin with severe pain
    \item Low blood pressure, fainting, or collapse
    \item A stiff neck with high fever and headache in a baby or child
    \item Any serious illness in a newborn baby under one month old
\end{itemize}

\section*{Diagnosis and Investigations}
Diagnosis usually combines clinical examination with targeted laboratory tests.

\begin{itemize}
    \item \textbf{Throat swab:} cultured to confirm group A streptococcus in suspected strep throat, or tested with a rapid antigen kit
    \item \textbf{Skin and wound swabs:} identify the bacteria in impetigo and cellulitis and their antibiotic sensitivities
    \item \textbf{Blood cultures:} taken when invasive infection is suspected, to grow the organism from the blood
    \item \textbf{Full blood count and inflammatory markers:} support the diagnosis of significant bacterial infection
    \item \textbf{Antibody tests:} checking for previous streptococcal infection when rheumatic fever or kidney disease is suspected
    \item \textbf{GBA screening in pregnancy:} vaginal and rectal swabs at around 36 weeks detect group B carriage so intrapartum antibiotics can be planned
    \item \textbf{Imaging:} chest X-ray for pneumonia, CT or ultrasound for deep collections of pus
    \item \textbf{Lumbar puncture:} cerebrospinal fluid is examined if meningitis is suspected
\end{itemize}

\section*{Management}
Treatment depends on the type and severity of infection.

\begin{itemize}
    \item \textbf{Antibiotics:} penicillin or amoxicillin is first-line for most group A infections, given for ten days for throat infections; erythromycin or azithromycin is used in penicillin allergy
    \item \textbf{Pain relief and fluids:} paracetamol or ibuprofen for fever and pain, with plenty of fluids
    \item \textbf{Skin care:} impetigo responds to topical antibiotics or, for widespread sores, a short course of oral antibiotics; sores should be kept clean and covered
    \item \textbf{Intrapartum antibiotics:} for group B carriers, intravenous penicillin during labour markedly reduces transmission to the baby
    \item \textbf{Newborn treatment:} babies with suspected group B infection are hospitalised for intravenous antibiotics and monitoring
    \item \textbf{Drainage of collections:} abscesses and deep pus collections may need surgical drainage
    \item \textbf{Intensive care:} invasive streptococcal disease and toxic shock require hospital care, intravenous antibiotics, and organ support
    \item \textbf{Prevention of recurrence:} people who have had rheumatic fever take long-term penicillin to prevent repeat infections that would damage their hearts further
    \item \textbf{Public health follow-up:} serious cases may be notified and contacts offered assessment
\end{itemize}

\section*{Complications}
Most streptococcal infections resolve completely, but complications are the reason treatment matters.

\begin{itemize}
    \item \textbf{Rheumatic fever and rheumatic heart disease:} immune damage to heart valves after throat infection, with arthritis, fever, and a characteristic movement disorder
    \item \textbf{Post-streptococcal glomerulonephritis:} kidney inflammation after throat or skin infection, causing blood in the urine, swelling, and high blood pressure
    \item \textbf{Neonatal sepsis and meningitis:} the most feared outcome of group B carriage, with high rates of death and long-term disability when it occurs
    \item \textbf{Invasive group A disease:} bloodstream infection, necrotising fasciitis, and toxic shock syndrome, which can be fatal despite treatment
    \item \textbf{Infective endocarditis:} valve infection from mouth streptococci, requiring prolonged antibiotics and sometimes surgery
    \item \textbf{Spread to adjacent sites:} middle-ear infection, sinusitis, and throat abscess following tonsillitis
    \item \textbf{Antibiotic-related problems:} allergic reactions, and overuse contributing to resistance
\end{itemize}

\section*{Prognosis and Outcomes}
With prompt and appropriate treatment, the outlook for the common streptococcal infections is excellent. Strep throat and impetigo resolve in nearly all cases, and rheumatic fever is preventable when throat infections are treated in time, which is why high-risk communities have dedicated prevention programs. Group B screening in pregnancy has reduced newborn disease dramatically, and babies who develop early infection generally recover well if treatment begins without delay. The serious forms, such as invasive group A disease and neonatal sepsis, remain dangerous, with significant mortality, but outcomes have improved with intensive care, and survival is usual when treatment starts early. The key messages are that most streptococcal infections are mild and easily treated, that prompt treatment of sore throats and skin sores prevents the long-term complications, and that anyone with features of invasive disease needs emergency care. Long-term survivors of rheumatic heart disease require lifelong follow-up, and penicillin prophylaxis protects against further damage.

\section*{Prevention and Self-Care}
\begin{itemize}
    \item Wash hands regularly, especially after coughing, sneezing, and touching the nose and mouth
    \item Cover coughs and sneezes, and stay home while unwell with a throat or skin infection
    \item Keep impetigo sores clean and covered, and avoid sharing towels and linen until they heal
    \item Attend the routine group B streptococcus screening offered in pregnancy and follow the birth plan agreed with the maternity team
    \item Take antibiotics exactly as prescribed and complete the full course
    \item Ensure children are up to date with the pneumococcal and other recommended vaccines
    \item Practise good wound hygiene for cuts and scrapes, and seek prompt treatment of any spreading redness
    \item In high-risk communities, have throat and skin infections treated early to prevent rheumatic fever
\end{itemize}

\section*{Sources and Further Reading}
The following organisations provide reliable information about streptococcal infections and their prevention.

\begin{itemize}
    \item NHS. \textit{Group A streptococcus: symptoms, treatment and when to seek urgent help.} https://www.nhs.uk/
    \item Mayo Clinic. \textit{Strep infections: causes, symptoms and treatment.} https://www.mayoclinic.org/
    \item healthdirect Australia. \textit{Streptococcal infections and scarlet fever.} https://www.healthdirect.gov.au/
    \item Australian Department of Health and Aged Care. \textit{Group B streptococcus and pregnancy; rheumatic fever control.} https://www.health.gov.au/
    \item World Health Organization (WHO). \textit{Group A streptococcus and rheumatic fever fact sheets.} https://www.who.int/
    \item RHDAustralia (Rheumatic Heart Disease Australia). \textit{ARF and RHD prevention guidelines.} https://www.rhdaustralia.org.au/
\end{itemize}

\clearpage
